\chapter{Findings} %This is an example title name

\begin{chapterabstract}

This is where you will insert your chapter abstract, which funcitons as a concise, self-contained summary of your chapter.

%Note that chapter abstracts are not required by KU. If your committee feels like you don't need an abstract for each chapter, simply don't include the \begin{chapterabstract} and \end{chapterabstract} commands in each chapter .tex file.

\end{chapterabstract}

%%%%%%%%%%%%%%%%%%%%%%%%%%%%%%%%%%%%%%%%%%%%%%%%%%%%%%%%%%%%%%%%%%%%%%%%%%%%%%%%%%%%%%%%%%
%%%%%%%%%%%%%%%%%%%%%%%%%%%%%%%%%%%%%%%%%%%%%%%%%%%%%%%%%%%%%%%%%%%%%%%%%%%%%%%%%%%%%%%%%%

\section{Name of Subsection}
Type out the content of your subsection here. Here's another citation \citep{Einstein:1933gu,Tho98w}. What follows is an example table:

\begin{table}[h]
\caption{Insert table caption here} % title of Table
\centering % used for centering table
\begin{tabular}{c c c c c} % centered columns (5 columns)
\hline\hline %inserts double horizontal lines
Set & Test 1 & Test 2 & Test 3 & Test 4 \\ [0.5ex] % inserts table
%heading
\hline % inserts single horizontal line
1 & 50 & 34 & 837 & 970 \\ % inserting body of the table
2 & 47 & 877 & 34 & 230 \\
3 & 5 & 31 & 25 & 415 \\
4 & 35 & 27 & 144 & 2356 \\
5 & 45 & 300 & 556 & 389 \\ [1ex] % [1ex] adds vertical space
\hline %inserts single line
\end{tabular}
\label{table_example_1:nonlin} % is used to refer to this table in the text
\end{table}

And just to be safe, here is another example of a table:

\begin{table}[h]
\caption{Insert table caption here} 
\centering 
\begin{tabular}{c c c c c c} 
\hline\hline 
Experiment & Batch A & Batch B & Batch C & Batch D & Batch E \\ [0.5ex] 
\hline 
5 & 7 & 5 & 3 & 7 & 2 \\
8 & 9 & 8 & 2 & 1 & 9 \\
3 & 2 & 1 & 9 & 6 & 6 \\
2 & 2 & 8 & 1 & 4 & 9 \\ [1ex] 
\hline 
\end{tabular}
\label{table_example_2:nonlin}
\end{table}

%The above table based on the one from "Creating Tables with LaTeX" (https://web.archive.org/web/20090306042457/http://www1.maths.leeds.ac.uk/latex/TableHelp1.pdf)

You can also refer to tables in the text like this: see tables \ref{table_example_1:nonlin} and \ref{table_example_2:nonlin}.

\section{Name of Subsection}
Here is another example of a chapter subsection. Delete or duplicate as necessary.

\subsection{Name of Third-Level Subsection}
Here is another example of a third-level subsection. Delete or duplicate as necessary.

\subsubsection{Name of Fourth-Level Subsection}
Here is another example of a fourth-level subsection. Delete or duplicate as necessary.